\documentclass[spanish, 12pt]{article}
\usepackage[T1]{fontenc} % https://tex.stackexchange.com/questions/664/why-should-i-use-usepackaget1fontenc
\usepackage[spanish,es-tabla]{babel} % Traducciones, y cambia cuadro por tabla
\usepackage{csquotes} % Mejor soporte de comillas en bibliografía
\usepackage{circuitikz} % Para armar circuitos directamente en LaTeX
\usepackage{amsmath} % Mejoras en fórmulas matemáticas
\usepackage{multirow} % Multples filas 
\usepackage{multicol} % Múltiples columnas
\usepackage[a4paper, portrait, margin=2.5cm]{geometry} % Ajusto márgenes y tamaño de página
\usepackage{booktabs} % Mejoras en tablas
\usepackage{graphicx} % Extended interface for graphics inclusion
\usepackage{float} % Mejor soporte de float para figuras
\usepackage{cancel} % Permite cancelar términos

% --- Para listar código
\usepackage{listings} 
\usepackage{xcolor}
% --- Para figuras y subfiguras
\usepackage{caption}
\usepackage{subcaption}
% --- Para el título
\usepackage[absolute,overlay]{textpos}
% --- Para hipervínculos
\usepackage{hyperref}
\hypersetup{
    hidelinks = true
}

\urlstyle{same}

% --- Colores para listados de código
\definecolor{codegreen}{rgb}{0,0.6,0}
\definecolor{codegray}{rgb}{0.5,0.5,0.5}
\definecolor{codepurple}{rgb}{0.58,0,0.82}
\definecolor{backcolour}{rgb}{0.95,0.95,0.92}

\lstdefinestyle{mystyle}{
    commentstyle=\color{codegreen},
    keywordstyle=\color{magenta},
    numberstyle=\tiny\color{codegray},
    stringstyle=\color{codepurple},
    basicstyle=\ttfamily\footnotesize,
    breakatwhitespace=false,
    breaklines=true,
    captionpos=b,
    keepspaces=true,
    numbers=left,
    numbersep=5pt,
    showspaces=false,
    showstringspaces=false,
    showtabs=false,
    tabsize=2
}

\lstset{style=mystyle}
% ---

% --- Configuración de bibiliografía
\usepackage[
backend=biber,
style=alphabetic,
sorting=ynt
]{biblatex}
\addbibresource{references.bib}
% ---

% --- Para informes de tipo pregunta/respuesta
\newcounter{question}
\setcounter{question}{0}
\newcommand{\question}[1]{
    \stepcounter{question}
    \vspace{0.2cm}
    \noindent\rule{\textwidth}{0.4mm}
    {\large\thequestion. #1\par}
    \noindent\rule{\textwidth}{0.4mm}
}
\newcommand{\answer}[1]{
    \vspace{0.2cm}
    #1
}

% --- Unidades
\usepackage{siunitx}
\sisetup{
    separate-uncertainty=true,
    locale = DE,
    round-mode = uncertainty,
    round-precision = 1
}

% --- Para TODOs y comentarios
\newcommand{\TODO}[1]{
    \begin{center}
        \colorbox{orange}{{\Huge TODO: }#1}
    \end{center}
}

\newcommand{\comment}[1]{
    \begin{center}
        \colorbox{yellow}{#1}
    \end{center}
}

% Graficos con gnuplot
\usepackage{tikzscale}
\usepackage{gnuplot-lua-tikz}

\begin{document}

    \newcommand{\materia}[1]{{\huge #1}\\\vspace{0.5cm}}
\newcommand{\titulo}[1]{{\Huge \textbf{#1}}\\\vspace{0.5cm}}
\newcommand{\subtitulo}[1]{{\huge #1}\\\vspace{0.5cm}}
\newcommand{\autor}[3]{#1\\{\small Mat. #2\\\href{mailto:#3}{#3}}\\\vspace{0.7cm}}
\newcommand{\fecha}[4]{{\small Fecha de #1: #4-#3-#2}\\} % YYYY-MM-DD
\newcommand{\departamento}[1]{{\small Departamento de #1\\Facultad de Ingeniería\\Universidad Nacional de Mar del Plata}}
\newenvironment{mytitlepage}
    {\pagenumbering{gobble}\begin{center}}
    {\end{center}\newpage\pagenumbering{arabic}}

\begin{mytitlepage}
    \includegraphics[width=4cm]{img/logo-fi.pdf}
   
    \begin{textblock*}{\paperwidth}(0cm, 7cm)
    \materia{Física Experimental}
    \titulo{Trabajo Práctico 6}
    \subtitulo{Péndulo de varilla 2}
    \end{textblock*}
    
    \begin{textblock*}{\paperwidth}(0cm, 12cm)
    \autor{Nazareno Montesoro}{13.668}{n.montesoro@gmail.com}
    \autor{Sofía Romero Ortueta}{13.778}{ssofi.romero@gmail.com}
    \end{textblock*}
    
    \begin{textblock*}{\paperwidth}(0cm, 25cm) 
    \fecha{realización}{28}{10}{2022}
    \fecha{entrega}{03}{11}{2022}
    \end{textblock*}
    
    \begin{textblock*}{\paperwidth}(0cm, 27cm) 
    \departamento{Física}
    \end{textblock*}
\end{mytitlepage}

    \tableofcontents

    \newpage
    \section{Resumen}

En esta práctica de laboratorio se montan circuitos electrónicos en una placa
de pruebas pasando por distintas configuraciones de amplificadores
operacionales, como la de amplificador inversor, amplificador no inversor
y la de seguidor de tensión. Asimismo se hace uso de un circuito integrado con
el cual se logra reducir el número de componentes necesarios para el montaje
en cuestión.

\TODO{¿Cambiar esa última oración? No sé si vienen operacionales ``por 
separado''}

\TODO{como para reducir el número de componentes usando un integrado.}

Por otro lado, se realiza un análisis de cada circuito junto a la
aplicación de simulaciones por computadora para posteriormente realizar
mediciones que verifiquen la eficacia de las ecuaciones que estudiamos
en el curso y de dichas simulaciones.

\TODO{Modificar la redacción.}


    \section{Introducción}

El período $T$ de un péndulo físico está dado por la ecuación
\ref{ec:intro:periodo}, donde $L$ es la longitud del péndulo y $g$ es la
constante de la gravedad:

\begin{equation}
    \label{ec:intro:periodo}
    T = 2\pi \sqrt{\frac{2L}{3g}}
\end{equation}

De la ecuación anterior puede despejarse $g$, de forma tal de poder calcular el
valor de esta constante (y su incerteza) midiendo la longitud y el tiempo que
tarda un péndulo en completar una oscilación:

\begin{equation}
    \label{ec:intro:gravedad}
    g = \frac{8}{3} \pi^2 \frac{L}{T^2}
\end{equation}

\begin{equation}
    \label{ec:intro:error-gravedad}
    \Delta g = \frac{8}{3} \pi^2 \left( \frac{\Delta L}{T^2} +
                                        \frac{2L \Delta T}{T^3}
                                    \right)
\end{equation}
donde $\Delta L$ y $\Delta T$ son las incertidumbres en las mediciones de la
longitud y el período del péndulo, respectivamente. 


    \section{Procedimiento}

Tome una varilla y mida su longitud con la cinta métrica. Luego utilice el 
soporte y aparte la varilla de su posición de equilibrio. Permita que el 
péndulo oscile con una amplitud pequeña y, utilizando un cronómetro, tome el
tiempo que demora \textbf{una única oscilación} (un período). Repita la 
operación y registre 200 valores del período que obtiene, manteniendo las
mismas condiciones en cada adquisición.

Una vez obtenidas todas las medidas para la primera varilla, repita el
procedimiento en una segunda varilla idéntica a la anterior y registreuna 
única medición del período.


    \section{Análisis de datos}

\begin{figure}[H]
    \centering
    \input{img/periodos.tikz}
    \caption{Períodos obtenidos}
\end{figure}

\begin{figure}[H]
    \centering
    \input{img/cuadrados.tikz}
    \caption{Mínimos cuadrados}
\end{figure}


    \newpage
    \section{Apéndice}
\label{sec:apendice}

\subsection{Fórmulas para el cálculo de estimadores en mínimos cuadrados}
\label{sec:apendice:formulas-estimadores-mc}

Sea $y_i = Bx_i + A$ una recta que ajusta una serie de $N$ datos mediante
cuadrados mínimos. La pendiente $B$ y la ordenada al origen $A$ se calculan
como:

\newcommand{\denominador}{N \sum {x_i}^2 - \left( \sum x_i \right)^2} 

\begin{equation}
    \label{ec:apendice:pendiente-mc}
    B = \frac{ N \sum x_i y_i - \sum x_i \sum y_i }
             { \denominador }
\end{equation}
\vspace{5mm}
\begin{equation}
    \label{ec:apendice:ordenada-mc}
    A = \frac{ \sum {x_i}^2 \sum y_i - \sum x_i \sum x_i y_i}
             { \denominador }
\end{equation}
\vspace{5mm}

A su vez, se puede calcular la desviación de $y_i$ como:

\begin{equation}
    \label{ec:apendice:desviacion-y-mc}
    \sigma_{y_i} = \sqrt{ \frac{1}{N-2} \sum 
                          \left( y_i - A - B x_i \right)^2 }
\end{equation}

\vspace{5mm}
La desviación de la pendiente y de la ordenada pueden calcularse como:

\begin{equation}
    \label{ec:apendice:desviacion-pendiente-mc}
    \sigma_B = \sigma_{y_i} \sqrt { \frac{N}{\denominador} }
\end{equation}

\vspace{5mm}
\begin{equation}
    \label{ec:apendice:desviacion-ordenada-mc}
    \sigma_A = \sigma_{y_i} \sqrt { \frac{ \sum {x_i}^2 }{ \denominador } }
\end{equation}

\vspace{5mm}
Finalmente se puede calcular el error en la pendiente y en la ordenada como:

\begin{equation}
    \label{ec:apendice:errores-mc}
    \Delta B = \frac{3\,\sigma_B}{\sqrt{N}}, \quad\quad
    \Delta A = \frac{3\,\sigma_A}{\sqrt{N}}
\end{equation}

    
    \newpage
    \section{Bibliografía}
    \printbibliography[heading=none]

\end{document}
